\textbf{结论名称：} 切线性质定理（切线垂直于过切点的半径）

\textbf{适用条件：} 直线与圆相切，且已知切点；或需由垂直关系反推切线。

\textbf{变量说明：} 设$\odot O$的圆心为$O$，半径为$r$，切点为$P$，切线为直线$l$，$P\in l$。

\textbf{核心公式：}
\[
\boxed{OP \perp l}
\]
即圆心与切点的连线（过切点的半径）垂直于切线。

\textbf{使用场景：} 连接圆心与切点构造直角三角形，常与勾股定理、垂径定理等结合求线段长度、角度或证明位置关系。

\textbf{简单例子：} 如图，$\odot O$的半径$r=3$，直线$l$切圆于点$P$，$Q$是$l$上一点且$PQ=4$。连接$OQ$，因$OP\perp l$，在$\mathrm{Rt}\triangle OPQ$中，由勾股定理得
\[
OQ = \sqrt{OP^{2}+PQ^{2}} = \sqrt{3^{2}+4^{2}} = 5.
\]
该例可直接手算验证垂直与直角三角形的关系。